% !PRINTED 552
\weberpage{552}
We suppose that the given values of $\cos m\varphi$ and $\sin m\varphi$ are less than 1 and that the sum of their squares is $= 1$.

We further adjoin $m$th roots of unity, with which we shall occupy ourselves more thoroughly in the next section, but which in any case depend on equations whose degree is lower than $m$. In order, however, to keep the field real, we shall adjoin, instead of the roots of unity themselves, the two real quantities
\weq{1}{\sin\frac{2\pi}{m}, \qquad \cos\frac{2\pi}{m}}

Let the field $\Omega$, then, consist of all rational numbers and of the rational functions of $\cos m\varphi$, $\sin m\varphi$, $\cos\frac{2\pi}{m}$, $\sin\frac{2\pi}{m}$.

The equation of degree $m$ on which $x = 2\cos\varphi$ depends we have set up in \S.\,137 in the form
\weq{2}{2\cos m\varphi = A_m(x),}
where $A_m(x)$ is a polynomial of degree $m$ in $x$; but there we have also found the equation
\weq{3}{\sin\varphi = \frac{\sin m\varphi}{B_m(x)}}
by which $\sin\varphi$ is expressed rationally in terms of $x$. In these formulae the roots of unity are not yet contained.

The $m$ roots of (2) now have the following meaning:
\[
x_0 = 2\cos\varphi, \qquad x_1 = 2\cos\left(\varphi + \frac{2\pi}{m}\right)\dots,
\]
\[
x_{m-1} = 2\cos\left(\varphi + \frac{2(m-1)\pi}{m}\right),
\]
and with the aid of (2) and (3), and upon adjunction of the quantities (1), each of them can be represented as a rational function of another, namely thus:
\[
x_1 = f(x_0), \qquad x_2 = f(x_1)\dots, \qquad x_0 = f(x_{m-1}).
\]

Any function of the $x_0, x_1 \dots x_{m-1}$ can therefore be represented as a rational function $F(x_0)$; under a cyclic permutation the following transformations occur:
\[
F(x_0) \text{ into } F(x_1), \ F(x_1) \text{ into } F(x_2) \dots F(x_{m-1}) \text{ into } F(x_0)
\]
For a \gesperrt{cyclic function}, therefore,
%CONT>>
